With the help of a rocket R, launched from a spaceship N, an asteroid A
must be destroyed, approaching the Earth threateningly, coming right in
the direction of the center of the Earth, as shown in Figure 1. The figure
shows the positions of the spacecraft and asteroid when the rocket is
launched. The spacecraft has constant Earth-relative velocity $\vec v_N$
along the spacecraft--asteroid line. The asteroid has constant Earth-relative
velocity $\vec v_A$, making an angle $\alpha$ with the spacecraft's direction.
% editorial: A duplicated and garbled phrase about v_N in the original has been consolidated here.
During the rocket's flight until impact, consider the Earth at rest relative
to the Sun, and the motions of the rocket R, spacecraft N, and asteroid A
rectilinear and uniform.
\ioaafig{2022_TH_Q12-f1.pdf}
We assume that the direction of absolute displacement of rocket R is
perpendicular to the direction of absolute movement of asteroid A.
Determine the angle $\beta$ at which rocket R was launched relative to the
direction of motion of spacecraft N.
